\documentclass[conference]{IEEEtran}
\IEEEoverridecommandlockouts
%immer rechts kapitel beginnen

%Default-einstellungen für das Dokument

\usepackage[T1]{fontenc}    %Schriftcodierung
\usepackage[utf8]{inputenc}    %Dateicodierung
\usepackage[ngerman]{babel} %Sprachpaket(Silbentrennung,überschriften, Abbildungen)
\usepackage{graphicx}
\usepackage{lmodern} %Vektorschrift
\usepackage{sansmath}
\usepackage{blindtext}
\usepackage{svg}
\usepackage{amsfonts}
\usepackage{amsmath}
\usepackage{svg}
\usepackage{caption}

%Literatur anpassen
\usepackage[
    backend=biber,
    style=numeric-verb,
    sorting=none % numeric-verb, numeric-comb
]{biblatex}
\addbibresource{ref.bib}

%Schriftanpassungen
\renewcommand{\familydefault}{\sfdefault} %serifenlose Schrift
\sansmath            %Aktivieren von Paket "Sansmath"

\usepackage[
    automark, %automatische Kolumne
    headsepline %Trennlinie in Kopfzeile
]{scrlayer-scrpage}
\pagestyle{scrheadings}
\clearscrheadfoot

\ohead[]{\headmark} %dynamische kolumne in kopfzeile
\ofoot[\pagemark]{\pagemark} %Seitenzahl in Fußzeile außen

%wichtige Pakete
\graphicspath{{bilder/}{}}  %Pfad zu den Bildern. mehrere pfade mit {{}{}{}} 

\title{Welche Rolle spielen künstliche Intelligenz und Machine Learning bei der Optimierung von Energieversorgungssystemen?}

\author{{\Large Nico Elsner}}




\begin{document}
	\ohead{\includegraphics[height=45pt]{bilder/thi.png.png}}
	\maketitle
	\begin{abstract}
		Die Optimierung von Energieversorgungssystemen ist angesichts der steigenden
Nachfrage nach effizienter, zuverlässiger und nachhaltiger Energieversorgung
von großer Bedeutung. In diesem Zusammenhang spielen künstliche Intelligenz
(KI) und maschinelles Lernen (ML) eine immer wichtigere Rolle. Dieser Artikel
untersucht die Rolle von KI und ML bei der Optimierung von
Energieversorgungssystemen. Es werden verschiedene Anwendungen von KI und ML
vorgestellt, darunter Prognose und Optimierung, Lastmanagement und
Nachfragesteuerung, Anlagenüberwachung und Wartung, Integration erneuerbarer
Energien sowie Energieeffizienz und Verbraucherempfehlungen. Darüber hinaus
werden die Herausforderungen diskutiert, die mit der Optimierung von
Energieversorgungssystemen verbunden sind, wie die Komplexität der Systeme, die
Volatilität erneuerbarer Energien, die Skalierbarkeit und Flexibilität, die
Sicherheit und Zuverlässigkeit sowie die wirtschaftlichen und regulatorischen
Aspekte. Die Erkenntnisse dieser Arbeit tragen dazu bei, das Verständnis für
die Rolle von KI und ML bei der Optimierung von Energieversorgungssystemen zu
vertiefen und zeigen Potenziale sowie Herausforderungen auf, die bei der
Einführung dieser Technologien berücksichtigt werden sollten.\\
	\end{abstract}
	%% Einleitung
	\section{Einleitung}
	\subsection{Hintergrund und Motivation des Themas}
	Hintergrund dieser Arbeit ist die in letzter Zeit extrem gestiegene Die
Energieversorgungssysteme stehen heutzutage vor einer Vielzahl von
Herausforderungen. Der steigende Energiebedarf, die begrenzte Verfügbarkeit
konventioneller Energieressourcen und die zunehmende Bedeutung erneuerbarer
Energien erfordern innovative Lösungen, um die Effizienz und Nachhaltigkeit der
Energieversorgung zu verbessern. In diesem Zusammenhang spielen künstliche
Intelligenz (KI) und maschinelles Lernen (ML) eine immer wichtigere Rolle.

Die Grundidee hinter dem Einsatz von KI und ML in der Optimierung von
Energieversorgungssystemen besteht darin, dass diese Technologien in der Lage
sind, große Mengen an Daten zu analysieren, Muster zu erkennen, Vorhersagen zu
treffen und adaptive Entscheidungen zu treffen. Dadurch können
Energieversorgungssysteme effizienter gestaltet, die Zuverlässigkeit erhöht und
Kosten gesenkt werden.

Ein zentrales Anwendungsgebiet von KI und ML liegt in der Prognose von
Energieerzeugung und -nachfrage. Durch den Einsatz von ML-Algorithmen können
präzisere Vorhersagen über den zukünftigen Energiebedarf gemacht werden, was
eine optimierte Planung der Energieerzeugung und -verteilung ermöglicht.
Darüber hinaus können ML-Modelle genutzt werden, um Lastmanagement-Strategien
zu entwickeln, um den Energieverbrauch zu optimieren und die Spitzenlasten zu
reduzieren.

Ein weiteres wichtiges Einsatzgebiet von KI und ML liegt in der Optimierung von
Stromnetzen und -verteilung. Hier können intelligente Algorithmen eingesetzt
werden, um den Energiefluss zu steuern, Engpässe zu identifizieren und die
Netzstabilität zu gewährleisten. ML kann auch bei der Erkennung von Fehlern und
Störungen in Energieversorgungssystemen helfen, indem Anomalien in den Daten
erkannt und potenzielle Probleme frühzeitig gemeldet werden.

Die Integration erneuerbarer Energien stellt ebenfalls eine große
Herausforderung dar, da sie von Natur aus volatil und schwer vorhersehbar sind.
Hier kommen KI und ML zum Einsatz, um Vorhersagemodelle für die Stromerzeugung
aus erneuerbaren Quellen zu entwickeln, um die Integration in das Gesamtsystem
zu optimieren. Dies ermöglicht eine bessere Ausnutzung erneuerbarer Energien
und eine verbesserte Integration mit konventionellen Energiequellen.

Neben den technischen Aspekten gibt es auch weitere wichtige Faktoren, die bei
der Optimierung von Energieversorgungssystemen berücksichtigt werden müssen.
Datenschutz und Sicherheit spielen eine entscheidende Rolle, da die
Verarbeitung großer Mengen an Energie- und Verbrauchsdaten eine Herausforderung
für den Schutz sensibler Informationen darstellt. Darüber hinaus stellen sich
auch ethische Fragen, wie beispielsweise der faire Zugang zu Energie oder die
Vermeidung von Diskriminierung bei der Nutzung von KI-gesteuerten
Energiesystemen.

In Anbetracht all dieser Aspekte ist es von großer Bedeutung, die Rolle von KI
und ML bei der Optimierung von Energieversorgungssystemen zu untersuchen. Durch
die Nutzung des Potenzials dieser Technologien können wir eine effizientere,
nachhaltigere und zuverlässigere Energieversorgung erreichen. In dieser
Seminararbeit werden wir die verschiedenen Einsatzbereiche von KI und ML in der
Energieversorgung untersuchen und ihre Auswirkungen sowie mögliche
Herausforderungen und zukünftige Entwicklungen analysieren.
	%\section{Zielsetzung und Fragestellung der Arbeit}
	%\ldots


	\section{Grundlagen der künstlichen Intelligenz und Machine Learning}
	\input{kapitel/Grundlagen der künstlichen Intelligenz und Machine Learning.tex}
	\subsection{Neuronale Netze}
	Künstliche neuronale Netze sind die Grundfunktionsweise von System, welche
künstliche Intelligenz nutzen. Diese Netze bestehen aus einer großen Menge an
Neuronen, welche, genauso wie die Neuronen im menschlichen Gehirn, eine
Vielzahl an Daten verarbeitet. In der folgenden Abbildung ist der Ablauf der
Arbeitsweise von künstlichen neuronalen Netzen exemplarisch dargestellt:

!!!!!!!!!! VEKTORGRAFIK EINBINDEN!!!!!!!!!!

Das Training eines neuronalen Netzes beginnt mit dem Bereitstellen der
Informationen, welches oft Signale oder Bildpixel sind. Zum Beispiel
Handschriften, Gesichter oder Daten von Wetterberichten. Darauffolgend wird der
Aufbau des neuronalen Netzes bestimmt, hinsichtlich der Anzahl an Hidden-Layers
und Neuronen, gefolgt von der Verarbeitung der Eingangsinformationen durch
Training mit verschiedenen Lernverfahren. Das Lernverfahren, mit dem geringsten
Fehler wird ausgewählt, aus welchem dann die Antwort des neuronalen Netzes in,
als Vektor kodierter Form ausgegeben wird.\cite{dorn2016programmieren}
	\subsection{Typen von Machine Learning Algorithmen}
	\begin{itemize}
    \item \textbf{Unsupervised Learning}:
          Beim unüberwachten Lernen gibt es keine Zielwerte, sondern das Ziel besteht
darin, Muster und Strukturen in den Daten zu entdecken. Das Modell lernt
eigenständig, wie es die Daten gruppieren oder Zusammenhänge finden kann.
Clustering und Dimensionsreduktion sind Beispiele für unüberwachtes Lernen.
    \item \textbf{Semi-supervised Learning}:
          Beim halbüberwachten Lernen werden Modelle mit einer Kombination aus gelabelten
und nicht gelabelten Daten trainiert. Somit agiert semi-supervised learning
zwischen unsupervised- und supervised learning. Motivation hinter der Art des
Lernens ist der Aufwand und die Kosten, des Anbieten von gelabelten Daten, bei
welchen Ein- und Ausgabewerte bekannt sind. \cite{lanquillon2019grundzuge} Das Modell nutzt die vorhandenen
markierten Daten, um Informationen zu generalisieren und die nicht markierten
Daten zu nutzen, um weitere Muster zu erkennen.
    \item \textbf{Supervised Learning}:
          Der Grundgedanke bei supervised Learning ist, dass die Beispiele (Ein- und
Ausgabewerte) bekannt sind. Die Aufgabe des Computers gleicht hier dem
Auswendiglernen, wobei der Rechner dies deutlich einfacher und schneller
schafft als Menschen. Hieraus ergibt sich eine Funktion, mit welcher bei
unbekannten Eingabewerten eine sinnvolle Ausgabe erzielt wird. Diese Funktion
wird in der künstlichen Intelligenz als Modell bezeichnet. Die Entwicklung des
Modells ist ein stetiger Prozess. Die Parameter des Modells werden während der
Lernphase immer wieder angepasst, um die Optimalen Werte hierfür zu finden. Der
besondere Vorteil am überwachten Lernen liegt darin, dass man die Ausgaben
kennt, und somit während des Lernvorgangs das Modell hinsichtlich der
Fehleranfälligkeit überprüfen kann. Bewertet wird das Modell dann jedoch mit
Daten, die noch nicht bekannt sind, wobei man möglichst versucht, zu vermeiden,
eine höhere Fehleranzahl zu haben, als in der Lernphase, genannt Overfitting.
Auch Underfitting stellt ein Problem dar, was meist der Fall ist, wenn die
Komplexität eines Modells zu niedrig ist.\cite{lanquillon2019grundzuge}\\ 
Man unterscheidet hierbei zwischen konkreten Ausgabewerten, sowie kontinuierliche
Werte. Diese beiden Wertearten bestimmen, ob es sich bei dem
supervised-learning Ansatz um eine Klassifikation oder eine Regression handelt.
Bei einem kontinuierlichen Ausgangstyp ist es eine Regression, bei einem
diskreten Ausgabewert spricht man von einer Klassifikation. Diese
Unterscheidung ist nur von den Ausgabetypen abhängig und nicht von den
Eingabewerten.\cite{matzka2021unuberwachtes}\\
    \item \textbf{Reinforcement Learning}:
          Reinforcement learning unterscheidet sich grundlegend von den bisher genannten
Lernarten. Während bei un-, semisupervised learning, sowie supervised learning
immer sehr viele Eingabedaten für das Lernsystem bereitgestellt werden müssen,
generiert reinforcement learning diese selbst.Ein sogenannter Agent interagiert
hierfür in einer bestimmten Umgebung. Das Ziel des Agenten besteht darin,
Aktionen auszuführen, um eine maximale Belohnung zu erhalten. Durch
Erfahrungen, Belohnungen und Bestrafungen passt der Agent sein Verhalten an, um
optimale Entscheidungen in der Umgebung zu treffen. Dieser Lernansatz ist
perfekt für Szenarien, in denen man keine, oder nur sehr schwierig
Trainingsdaten bereitstellen kann. \cite{lanquillon2019grundzuge}

\begin{figure}
    \centering
    \includesvg[width=0.4\textwidth]{bilder/reinforcementlearning.svg}
    \caption[width=0.4\textwidth]{Ablauf bei Reinforcement Learning}
    \label{fig:disreinforcementlearning}
\end{figure}
\end{itemize}

	\section{Optimierung von Energieversorgungssystemen}
	\input{kapitel/Optimierung von Energieversorgungssystemen.tex}
	\subsection{Überblick über Energieversorgungssystemen}
	\input{kapitel/Überblick von Energieversorgungssystemen.tex}
	\subsection{Herausforderungen bei der Optimierung}
	Die Optimierung von Energieversorgungssystemen steht vor einer Reihe von
Herausforderungen, die es zu bewältigen gilt. Eine der zentralen
Herausforderungen besteht in der Komplexität dieser Systeme.
Energieversorgungssysteme bestehen aus einer Vielzahl von Komponenten und
Akteuren, wie beispielsweise Energieerzeugung, -übertragung, -verteilung und
-verbrauch. Die Koordination und Optimierung all dieser Komponenten erfordert
eine umfassende Analyse und Steuerung, um effiziente Ergebnisse zu erzielen.

Ein weiterer wichtiger Aspekt ist die Volatilität erneuerbarer Energien. Die
Integration von erneuerbaren Energien wie Solarenergie und Windenergie in das
Energieversorgungssystem stellt eine Herausforderung dar. Diese Energiequellen
unterliegen starken Schwankungen aufgrund von Wetterbedingungen und weisen
volatilere Erzeugungsmuster auf als konventionelle Energieressourcen.

Darüber hinaus stehen Energieversorgungssysteme vor Herausforderungen im
Zusammenhang mit der Skalierbarkeit und Flexibilität. Da die Energienachfrage
schwankt und sich die Energieerzeugungsmuster ändern, müssen die Systeme in der
Lage sein, sich an diese Veränderungen anzupassen und eine zuverlässige
Energieversorgung sicherzustellen.

Die Sicherheit und Zuverlässigkeit des Energieversorgungssystems ist eine
weitere Herausforderung. Da das Energienetz ein attraktives Ziel für
Cyberangriffe darstellt, müssen Maßnahmen ergriffen werden, um die Systeme vor
potenziellen Bedrohungen zu schützen und die Verfügbarkeit der
Energieversorgung zu gewährleisten.

Nicht zuletzt sind auch wirtschaftliche und regulatorische Herausforderungen zu
berücksichtigen. Die Optimierung von Energieversorgungssystemen erfordert
Investitionen in Infrastruktur, Technologien und Forschung. Zudem müssen
regulatorische Rahmenbedingungen geschaffen werden, die die Integration neuer
Technologien und Geschäftsmodelle ermöglichen und den Übergang zu einer
nachhaltigen Energieversorgung unterstützen.

Insgesamt stehen die Optimierung von Energieversorgungssystemen vor
verschiedenen Herausforderungen, die von der Komplexität der Systeme über die
Integration erneuerbarer Energien bis hin zu Fragen der Sicherheit,
Skalierbarkeit und Wirtschaftlichkeit reichen. Die Bewältigung dieser
Herausforderungen erfordert einen interdisziplinären Ansatz, der
technologische, wirtschaftliche, regulatorische und soziale Aspekte
berücksichtigt.

	\subsection{Allgemeine Möglichkeiten der Optimierung}
	\input{kapitel/Allgemeine Möglichkeiten der Optimierung.tex}
	

	\section{Rolle von künstlicher Intelligenz und Machine Learning bei der Optimierung von Energieversorgungssystemen}
	\input{kapitel/Rolle von künstlicher Intelligenz und Machine Learning bei der Optimierung von Energieversorgungssystemen.tex}
	\subsection{Potenzial und Vorteile von künstlicher Intelligenz und Machine Learning}
	\input{kapitel/Potenzial und Vorteile von künstlicher Intelligenz und Machine Learning.tex}
	\subsection{Beispiele und Anwendungen von künstlicher Intelligenz und Machine Learning in der Energieversorgung}
	\input{kapitel/Beispiele und Anwendungen von künstlicher Intelligenz und Machine Learning in der Energieversorgung.tex}


	\section{Kritische Betrachtung und Ausblick}
	\input{kapitel/Kritische Betrachtung und Ausblick.tex}
	\subsection{Grenzen und Risiken von künstlicher Intelligenz}
	\input{kapitel/Grenzen und Risiken von künstlicher Intelligenz.tex}
	\subsection{Perspektiven}
	Die Rolle von künstlicher Intelligenz und maschinellem Lernen bei der
Optimierung von Energieversorgungssystemen bietet vielversprechende
Perspektiven für die Zukunft. Das die Softwareindustrie, welche die KI-Systeme
maßgeblich mitprägt, immer schneller am wachsen ist zeigt
Abbildung(\ref{fig:wachstumsberufe}). Unter allen Berufen, ist der Informatiker
mit großem Abstand der schnellst Wachsende.
\begin{figure}[!h]
    \centering
    \includesvg[width=1\columnwidth]{bilder/Wachstumsberufe Deutschland_2.svg}
    \caption[width=0.4\textwidth]{Ablauf bei Reinforcement Learning}
    \label{fig:wachstumsberufe}
\end{figure}
\cite{eichhorst2015zukunft}

Mit dem weiteren Fortschritt von KI- und ML-Technologien eröffnen sich neue
Möglichkeiten zur Effizienzsteigerung, Kostensenkung und Verbesserung der
Nachhaltigkeit in der Energieversorgung.

Eine vielversprechende Perspektive besteht in der Integration von KI und ML in
Smart Grids und intelligenten Energienetzen. Durch den Einsatz
fortschrittlicher Algorithmen und Modelle können Energieflüsse optimiert,
Spitzenlasten reduziert und erneuerbare Energien effizienter in das
Energiesystem integriert werden. Dies ermöglicht eine bessere Ausbalancierung
von Angebot und Nachfrage und trägt zu einer stabilen und nachhaltigen
Energieversorgung bei.

Des Weiteren eröffnet die Kombination von KI und ML mit dem Internet der Dinge
(IoT) und dem Einsatz von Sensoren neue Perspektiven für die Energieeffizienz.
Intelligente Geräte und Systeme können Energieverbrauchsdaten erfassen und
analysieren, um personalisierte Empfehlungen zur Reduzierung des
Energieverbrauchs zu geben. Dies ermöglicht es Verbrauchern, bewusstere
Entscheidungen zu treffen und ihren Energieverbrauch zu optimieren.

Ein weiteres vielversprechendes Feld ist die Entwicklung autonomer
Energieversorgungssysteme. Durch den Einsatz von KI und ML können autonome
Energieerzeugungs- und Speicherlösungen entwickelt werden, die den Bedarf an
menschlichem Eingriff minimieren. Diese Systeme können den Energieverbrauch und
die -erzeugung in Echtzeit überwachen, Vorhersagen treffen und automatisch
optimale Entscheidungen treffen, um die Effizienz und Zuverlässigkeit der
Energieversorgung zu maximieren.

Darüber hinaus bieten KI und ML Möglichkeiten zur besseren Vorhersage und
Bewältigung von Energienotfällen. Durch die Analyse von Echtzeitdaten können
frühzeitige Warnungen und reaktionsschnelle Maßnahmen bei Störungen im
Energiesystem ermöglicht werden. Dies unterstützt die Sicherheit und Resilienz
der Energieinfrastruktur.

Insgesamt eröffnen die Fortschritte in KI und ML spannende Perspektiven für die
Energieversorgungssysteme der Zukunft. Mit einer kontinuierlichen Forschung und
Entwicklung, der Verfügbarkeit hochwertiger Daten und einer sorgfältigen
Integration dieser Technologien können wir eine effizientere, nachhaltigere und
zuverlässigere Energieversorgung erreichen.


	\section{Fazit und Ausblick}
	Die Rolle von künstlicher Intelligenz (KI) und maschinellem Lernen (ML) bei der
Optimierung von Energieversorgungssystemen ist von großer Bedeutung. Die
Anwendung von KI und ML bietet vielfältige Möglichkeiten zur
Effizienzsteigerung, Kostenoptimierung, Integration erneuerbarer Energien und
Verbesserung des Verbraucherverhaltens. Durch den Einsatz von KI können
Energieversorgungsunternehmen fundierte Entscheidungen treffen, Prognosen
erstellen, den Energiefluss optimieren und die Zuverlässigkeit der
Energieversorgung verbessern.

Jedoch sind bei der Nutzung von KI in der Energieversorgung auch
Herausforderungen zu berücksichtigen. Grenzen wie die Qualität und
Verfügbarkeit von Daten, ethische Überlegungen, potenzielle
Fehlinterpretationen oder Fehlentscheidungen von KI-Modellen sowie
Sicherheitsrisiken erfordern eine sorgfältige Planung, Überwachung und
Implementierung.

Ausblick: Für die Zukunft bieten sich auf dem Gebiet der künstlichen
Intelligenz und des maschinellen Lernens in der Energieversorgung zahlreiche
Möglichkeiten. Fortschritte in der Datenverfügbarkeit und -qualität werden die
Genauigkeit und Leistungsfähigkeit von KI-Modellen verbessern. Eine verstärkte
Berücksichtigung ethischer Prinzipien bei der Entwicklung und Implementierung
von KI-Systemen wird zu einer fairen und transparenten Nutzung beitragen.

Darüber hinaus könnten Fortschritte im Bereich des Verständnis von
KI-Modellen dazu beitragen, Vertrauen und Akzeptanz in deren Entscheidungen zu
fördern. Die Integration von KI-Systemen mit dem Internet der Dinge (IoT) und
anderen fortschrittlichen Technologien ermöglicht eine noch intelligentere und
effizientere Energieversorgung.

Ein wichtiger Aspekt für den Ausblick ist die kontinuierliche Forschung und
Entwicklung auf dem Gebiet der künstlichen Intelligenz und des maschinellen
Lernens in der Energieversorgung. Neue Algorithmen, Modelle und Techniken
werden entwickelt, um die Herausforderungen anzugehen und die
Leistungsfähigkeit von KI-Systemen weiter zu verbessern.

Insgesamt bietet die Kombination von künstlicher Intelligenz und maschinellem
Lernen enorme Potenziale für die Optimierung von Energieversorgungssystemen.
Mit einer sorgfältigen Betrachtung der Grenzen, Risiken und Chancen können wir
eine nachhaltigere, effizientere und zuverlässigere Energieversorgung der
Zukunft erreichen.

	\printbibliography
\end{document}